% Generated draft: requires mathematical and source review.
\documentclass[11pt]{article}
\usepackage[margin=25mm]{geometry}
\usepackage{fontspec}
\setmainfont{Latin Modern Roman}
\usepackage{amsmath,amssymb,mathtools,unicode-math}
\setmathfont{Latin Modern Math}
\usepackage{enumitem,needspace,xurl,hyperref}
\setlength{\parindent}{0pt}
\setlength{\parskip}{5pt}
\setlength{\emergencystretch}{3em}
\newcommand{\N}{\mathbb{N}}
\newcommand{\Z}{\mathbb{Z}}
\newcommand{\Q}{\mathbb{Q}}
\newcommand{\R}{\mathbb{R}}
\newcommand{\C}{\mathbb{C}}
\newcommand{\F}{\mathbb{F}}
\newcommand{\A}{\mathbb{A}}
\newcommand{\PP}{\mathbb{P}}
\newcommand{\E}{\mathbb{E}}
\newcommand{\eps}{\varepsilon}
\newcommand{\dd}{\,\mathrm{d}}
\begin{document}
\section*{UnsolvedMath ID: 20002088}
\textbf{Source problem code:} catalogue entry
\subsection{Definitions and mathematical statement}
Eliashberg's Question 1.30 asks which topology or smooth structure of a manifold \(M\) is remembered by the canonical symplectic cotangent bundle \(T^*M\) or contact unit cosphere bundle \(ST^*M\). In particular, can homeomorphic but nondiffeomorphic bases, such as exotic spheres, be distinguished by nonsymplectomorphic full cotangent bundles or noncontactomorphic cosphere bundles? The question also asks whether the natural assignment of knot conormals in \(ST^*S^3\cong S^3\times S^2\) is injective. The source does not specify all equivalence, marking, or orientation conventions.

The aid's full-cotangent transfer lemma assumes closed smooth bases \(M_1,M_2\), Liouville forms \(\lambda_i\), and a symplectomorphism
\(\Phi:(T^*M_1,\dd\lambda_1)\to(T^*M_2,\dd\lambda_2)\).
Let \(i_1\) be the zero section and \(\pi_1\) the projection. Put
\[
 \beta=\Phi^*\lambda_2-\lambda_1,\qquad \alpha=i_1^*\beta.
\]
Then \(\beta\) and \(\alpha\) are closed, and the zero-section deformation retraction gives \(\beta=\pi_1^*\alpha+\dd h\) for a smooth function \(h\). Closed-form fiber translation
\[
 \tau_{-\alpha}(q,p)=(q,p-\alpha_q)
\]
is symplectic, and \(F=\Phi\circ\tau_{-\alpha}\) satisfies
\[
 F^*\lambda_2-\lambda_1=\dd(h\circ\tau_{-\alpha}).
\]
Thus every such full-cotangent symplectomorphism can be exactified by precomposition. The image \(F(i_1(M_1))\) is a compact exact Lagrangian copy of \(M_1\), and its inclusion in \(T^*M_2\) is a homotopy equivalence. This gives a transfer principle, not a proof that the bases are diffeomorphic or a resolution for all exotic spheres.

For the contact version choose unit cosphere forms \(\eta_i=\lambda_i|_{ST^*M_i}\). A coorientation-preserving contactomorphism \(\psi\) has
\(\psi^*\eta_2=f\eta_1\) for \(f>0\). Under the radial identification
\(T^*M_i\setminus0\cong\R\times ST^*M_i\) with Liouville form \(e^s\eta_i\), the aid gives the strict lift
\[
 \Psi(s,y)=(s-\log f(y),\psi(y)),\qquad
 \Psi^*(e^s\eta_2)=e^s\eta_1.
\]
This is defined only on punctured cotangent bundles. It need not extend smoothly over the zero sections; no exact-Lagrangian conclusion about the bases follows automatically from it. No assertion for coorientation-reversing maps is included.

For a smooth knot \(K\subset S^3\), its unit conormal is
\[
 \Lambda_K=\{(q,p)\in ST^*S^3:q\in K,\ p|_{T_qK}=0\},
\]
a Legendrian torus. A usual interpretation asks whether Legendrian isotopy of \(\Lambda_{K_0}\) and \(\Lambda_{K_1}\) forces isotopy of the knots, but the excerpt leaves oriented/unoriented knots and the target equivalence relation unstated. Legendrian isotopy, arbitrary ambient contactomorphism, and contact equivalence preserving additional fiber/peripheral markings must be distinguished. The source calls Ng's work the first result in this direction; this is historical attribution, not an updated completeness claim. The aid says known exact-Lagrangian and knot-conormal results do not settle the full unmarked contact question through the transfer statements alone.
\subsection{Short English statement}
The source asks how much of a base manifold or knot is remembered by cotangent and cosphere contact geometry. The aid exactifies a symplectomorphism of full cotangent bundles and obtains an exact Lagrangian homotopy-equivalent copy of the base. A cosphere contactomorphism lifts only off the zero sections, leaving an extension gap and unresolved equivalence conventions for the knot question.
\subsection{Sources}
\begin{itemize}
\item Local UnsolvedMath registry, permanent ID 20002088, exact source code catalogue entry; canonical \texttt{aim-geometry-notes.json}, \texttt{notes[425]}, tag \texttt{question}; AIM Workshop Problem Lists, Geometry, \emph{Holomorphic curves in contact geometry}, Question 1.30 attributed to Eliashberg. Created and updated 14 August 2026.
\item Supplied source URL \url{https://aimath.org/WWN/contactgeom2/contactgeom2.pdf}. The verification repairs split numbering 1.30, line hyphenation, and the smooth cosphere-product isomorphism sign from the official workshop rendering. No external source was accessed here.
\item Ng is named for early knot-conormal work. The aid mentions known exact-Lagrangian and conormal results but supplies no full publication entries, theorem numbers, or separate URLs.
\item Attempt 1: \texttt{partial\_result}, problem status at run \texttt{partially\_solved}, rigor label \texttt{proved}, novelty confidence \texttt{low}; classification \texttt{PARTIAL-PROGRESS}, registry status \texttt{partially\_solved}. Candidate: full-cotangent exactification and strict punctured contact lift. These are dated supplied assessments, not a full rigidity theorem.
\item Ulam AI / UnsolvedMath Contributors; CC BY 4.0. Machine-generated research aid, not peer reviewed; independently verify before citation.
\end{itemize}
\end{document}
